% ============================================================
% 讲义 — 人教版九年级上册 · 第21章 一元二次方程
% 编译命令：xelatex -interaction=nonstopmode 第21章_一元二次方程.tex
% ============================================================
\documentclass[12pt, a4paper]{ctexart}
\usepackage[margin=2.2cm]{geometry}
\usepackage{amsmath, amssymb}
\usepackage{enumitem}
\usepackage{tcolorbox}
\usepackage{booktabs}
\usepackage{array}
\usepackage{fancyhdr}
\usepackage{xcolor}
\usepackage{tikz}
\usepackage{hyperref}
\usepackage{graphicx}
\usepackage{multirow}
\usepackage{pgfplots}
\pgfplotsset{compat=1.18}

% 颜色定义
\definecolor{mainblue}{RGB}{41, 98, 168}
\definecolor{lightblue}{RGB}{230, 240, 255}
\definecolor{tipgreen}{RGB}{34, 139, 87}
\definecolor{lightgreen}{RGB}{235, 250, 240}
\definecolor{warnred}{RGB}{200, 50, 50}
\definecolor{lightred}{RGB}{255, 240, 240}
\definecolor{examorange}{RGB}{210, 120, 20}
\definecolor{lightorange}{RGB}{255, 248, 235}
\definecolor{apppurple}{RGB}{120, 60, 170}
\definecolor{lightpurple}{RGB}{245, 238, 255}

% 页眉页脚
\pagestyle{fancy}
\fancyhf{}
\fancyhead[L]{\small\textcolor{mainblue}{\textbf{人教版九年级上册 · 第21章 一元二次方程}}}
\fancyhead[R]{\small\textcolor{mainblue}{家教讲义}}
\fancyfoot[C]{\small —\ \thepage\ —}
\renewcommand{\headrulewidth}{0.6pt}
\renewcommand{\headrule}{\hbox to\headwidth{\color{mainblue}\leaders\hrule height \headrulewidth\hfill}}
\renewcommand{\footrulewidth}{0pt}

% tcolorbox 样式
\newtcolorbox{knowledgebox}[1][]{%
  colback=lightblue, colframe=mainblue,
  fonttitle=\bfseries\large, title=#1,
  arc=3mm, boxrule=0.8pt,
  left=8pt, right=8pt, top=6pt, bottom=6pt
}

\newtcolorbox{tipbox}[1][]{%
  colback=lightgreen, colframe=tipgreen,
  fonttitle=\bfseries, title=#1,
  arc=2mm, boxrule=0.6pt,
  left=6pt, right=6pt, top=4pt, bottom=4pt
}

\newtcolorbox{warnbox}[1][]{%
  colback=lightred, colframe=warnred,
  fonttitle=\bfseries, title=#1,
  arc=2mm, boxrule=0.6pt,
  left=6pt, right=6pt, top=4pt, bottom=4pt
}

\newtcolorbox{examplebox}[1][]{%
  colback=lightorange, colframe=examorange,
  fonttitle=\bfseries, title=#1,
  arc=3mm, boxrule=0.8pt,
  left=8pt, right=8pt, top=6pt, bottom=6pt
}

\newtcolorbox{appbox}[1][]{%
  colback=lightpurple, colframe=apppurple,
  fonttitle=\bfseries, title=#1,
  arc=3mm, boxrule=0.8pt,
  left=8pt, right=8pt, top=6pt, bottom=6pt
}

\setCJKmainfont{Hiragino Sans GB}
\setCJKsansfont{Heiti SC}

\begin{document}

% ============ 标题页 ============
\begin{center}
{\Huge\bfseries\textcolor{mainblue}{第21章\quad 一元二次方程}}\\[12pt]
{\Large\textcolor{gray}{家\quad 教\quad 讲\quad 义}}\\[8pt]
\rule{0.6\textwidth}{1pt}\\[8pt]
{\normalsize 适用：人教版九年级上册\quad |\quad 课时：6$\sim$8课时}
\end{center}
\vspace{10pt}

% ============================================================
% 第一部分：知识点讲解
% ============================================================
\section*{\textcolor{mainblue}{一、知识点讲解}}

\subsection*{21.1\quad 一元二次方程的概念}

\begin{knowledgebox}[一元二次方程的定义]
只含有一个未知数，且未知数的最高次数是$2$的整式方程叫做一元二次方程。
\medskip
一元二次方程的一般形式为：
\[
\boxed{ax^2 + bx + c = 0 \quad (a \neq 0)}
\]
其中 $a$ 叫做二次项系数，$b$ 叫做一次项系数，$c$ 叫做常数项。
\end{knowledgebox}

\begin{warnbox}[易错提醒]
\begin{enumerate}[leftmargin=*]
  \item 必须满足 $a \neq 0$，否则就不是一元二次方程。
  \item 一般形式中各项必须在等号左边，右边为 $0$。
  \item 判断一元二次方程时，需要先化简整理，再看未知数的最高次数。
\end{enumerate}
\end{warnbox}

\subsection*{21.2\quad 一元二次方程的解法}

\begin{knowledgebox}[解法一：直接开平方法]
如果方程能化为 $x^2 = p$ 或 $(mx+n)^2 = p$ 的形式（其中 $p \geqslant 0$），则可以直接开平方求解。
\[
x^2 = p \implies x = \pm\sqrt{p}
\]
\[
(mx+n)^2 = p \implies mx+n = \pm\sqrt{p} \implies x = \frac{-n \pm \sqrt{p}}{m}
\]
\end{knowledgebox}

\begin{knowledgebox}[解法二：配方法]
将方程 $ax^2+bx+c=0$ 化为完全平方的形式求解。
\medskip
步骤：
\begin{enumerate}[leftmargin=*]
  \item 若 $a \neq 1$，方程两边同除以 $a$，化为 $x^2 + \dfrac{b}{a}x + \dfrac{c}{a} = 0$。
  \item 将常数项移到右边：$x^2 + \dfrac{b}{a}x = -\dfrac{c}{a}$。
  \item 两边同时加上 $\left(\dfrac{b}{2a}\right)^2$，使左边成为完全平方式。
  \item 左边写成 $\left(x + \dfrac{b}{2a}\right)^2$ 的形式，再用开平方法求解。
\end{enumerate}
\[
\boxed{x^2 + \frac{b}{a}x + \left(\frac{b}{2a}\right)^2 = -\frac{c}{a} + \left(\frac{b}{2a}\right)^2}
\]
\end{knowledgebox}

\begin{knowledgebox}[解法三：公式法]
对于一般形式 $ax^2+bx+c=0$（$a\neq 0$），当 $\Delta = b^2-4ac \geqslant 0$ 时，方程的两个根为：
\[
\boxed{x = \frac{-b \pm \sqrt{b^2-4ac}}{2a}}
\]
这就是一元二次方程的\textbf{求根公式}，它适用于所有的一元二次方程。
\end{knowledgebox}

\begin{knowledgebox}[解法四：因式分解法]
如果一元二次方程的左边能分解为两个一次因式的乘积，即：
\[
ax^2+bx+c = a(x-x_1)(x-x_2) = 0
\]
则由 $a \neq 0$ 可得 $x - x_1 = 0$ 或 $x - x_2 = 0$，从而 $x = x_1$ 或 $x = x_2$。
\medskip
常用的分解方法：
\begin{enumerate}[leftmargin=*]
  \item \textbf{提公因式法}：如 $x^2-3x = 0 \implies x(x-3)=0$。
  \item \textbf{平方差公式}：如 $x^2-4=0 \implies (x+2)(x-2)=0$。
  \item \textbf{十字相乘法}：如 $x^2-5x+6=0 \implies (x-2)(x-3)=0$。
\end{enumerate}
\end{knowledgebox}

\begin{tipbox}[解法选择策略]
\begin{enumerate}[leftmargin=*]
  \item 方程形如 $ax^2+c=0$（缺一次项）$\longrightarrow$ 直接开平方法。
  \item 方程一边能因式分解 $\longrightarrow$ 因式分解法（最快）。
  \item 系数简单且需要推导公式 $\longrightarrow$ 配方法。
  \item 通用方法（万能） $\longrightarrow$ 公式法。
\end{enumerate}
\end{tipbox}

\subsection*{21.3\quad 判别式}

\begin{knowledgebox}[判别式 $\Delta$]
对于一元二次方程 $ax^2+bx+c=0$（$a \neq 0$），其判别式为：
\[
\boxed{\Delta = b^2 - 4ac}
\]
判别式与根的关系：
\begin{itemize}[leftmargin=*]
  \item $\Delta > 0$：方程有\textbf{两个不相等的实数根}。
  \item $\Delta = 0$：方程有\textbf{两个相等的实数根}（即一个重根）。
  \item $\Delta < 0$：方程\textbf{没有实数根}（有两个共轭复数根）。
\end{itemize}
\end{knowledgebox}

\begin{warnbox}[易错提醒]
判别式只在 $a \neq 0$ 的前提下才有意义。若题目含有参数，需要先讨论 $a$ 是否为 $0$，再计算 $\Delta$。
\end{warnbox}

\subsection*{21.4\quad 韦达定理（根与系数的关系）}

\begin{knowledgebox}[韦达定理]
若一元二次方程 $ax^2+bx+c=0$（$a\neq 0$）的两个实数根为 $x_1$、$x_2$，则：
\[
\boxed{x_1 + x_2 = -\frac{b}{a}, \qquad x_1 \cdot x_2 = \frac{c}{a}}
\]
特别地，当 $a=1$ 时，方程 $x^2+px+q=0$ 的两根满足：
\[
x_1 + x_2 = -p, \qquad x_1 \cdot x_2 = q
\]
\end{knowledgebox}

\begin{tipbox}[韦达定理的常见变形]
已知 $x_1+x_2 = S$，$x_1 \cdot x_2 = P$，则：
\begin{itemize}[leftmargin=*]
  \item $x_1^2 + x_2^2 = (x_1+x_2)^2 - 2x_1 x_2 = S^2 - 2P$
  \item $|x_1 - x_2| = \sqrt{(x_1+x_2)^2 - 4x_1 x_2} = \sqrt{S^2 - 4P}$
  \item $\dfrac{1}{x_1} + \dfrac{1}{x_2} = \dfrac{x_1+x_2}{x_1 x_2} = \dfrac{S}{P}$\quad（$P \neq 0$）
  \item $(x_1-x_2)^2 = (x_1+x_2)^2 - 4x_1 x_2 = S^2 - 4P$
\end{itemize}
\end{tipbox}

% ============================================================
% 第二部分：例题精讲
% ============================================================
\section*{\textcolor{examorange}{二、例题精讲}}

\begin{examplebox}[【例1】用直接开平方法解方程]
解方程：$(2x-3)^2 = 25$
\end{examplebox}

\textbf{解：}

由 $(2x-3)^2 = 25$，两边开平方得：
\[
2x - 3 = \pm 5
\]
当 $2x-3 = 5$ 时，$2x = 8$，$x = 4$；

当 $2x-3 = -5$ 时，$2x = -2$，$x = -1$。

所以方程的解为 $\boxed{x_1 = 4,\; x_2 = -1}$。

\begin{examplebox}[【例2】用配方法解方程]
用配方法解方程：$2x^2 - 4x - 6 = 0$
\end{examplebox}

\textbf{解：}

两边同除以 $2$：$x^2 - 2x - 3 = 0$。

移项：$x^2 - 2x = 3$。

配方：$x^2 - 2x + 1 = 3 + 1$，即 $(x-1)^2 = 4$。

开平方：$x - 1 = \pm 2$，所以 $\boxed{x_1 = 3,\; x_2 = -1}$。

\begin{examplebox}[【例3】用公式法解方程]
用求根公式解方程：$3x^2 + 5x - 2 = 0$
\end{examplebox}

\textbf{解：}

这里 $a=3$，$b=5$，$c=-2$。

计算判别式：$\Delta = b^2 - 4ac = 25 - 4 \times 3 \times (-2) = 25 + 24 = 49 > 0$。

代入求根公式：
\[
x = \frac{-b \pm \sqrt{\Delta}}{2a} = \frac{-5 \pm \sqrt{49}}{2 \times 3} = \frac{-5 \pm 7}{6}
\]
所以 $\boxed{x_1 = \dfrac{1}{3},\; x_2 = -2}$。

\begin{examplebox}[【例4】用因式分解法解方程]
解方程：$x^2 - 5x + 6 = 0$
\end{examplebox}

\textbf{解：}

用十字相乘法分解：寻找两个数，乘积为 $6$，和为 $-5$。

这两个数为 $-2$ 和 $-3$，因此：
\[
x^2 - 5x + 6 = (x-2)(x-3) = 0
\]
所以 $x-2=0$ 或 $x-3=0$，即 $\boxed{x_1 = 2,\; x_2 = 3}$。

\begin{tipbox}[十字相乘法口诀]
对于 $x^2+px+q=0$，寻找两个数 $m$、$n$，使得 $m+n=p$，$mn=q$，则 $x^2+px+q=(x+m)(x+n)$。
\end{tipbox}

\begin{examplebox}[【例5】判别式的应用]
已知关于 $x$ 的方程 $kx^2-2x+1=0$ 有两个不相等的实数根，求 $k$ 的取值范围。
\end{examplebox}

\textbf{解：}

因为方程有两个不相等的实数根，所以方程必须是一元二次方程，即 $k \neq 0$。

且判别式 $\Delta > 0$：
\[
\Delta = (-2)^2 - 4k \times 1 = 4 - 4k > 0
\]
解得 $k < 1$。

又因为 $k \neq 0$，所以 $k$ 的取值范围为 $\boxed{k < 1 \text{ 且 } k \neq 0}$。

\begin{warnbox}[易错提醒]
本题容易遗漏 $k \neq 0$ 的条件。若 $k=0$，方程变为 $-2x+1=0$，这是一元一次方程，不是一元二次方程，不能用判别式讨论。
\end{warnbox}

\begin{examplebox}[【例6】韦达定理的综合应用]
已知方程 $x^2-3x+1=0$ 的两个根为 $x_1$、$x_2$，不解方程，求下列各式的值：

$(1)\; x_1^2 + x_2^2 \qquad (2)\; \dfrac{1}{x_1} + \dfrac{1}{x_2} \qquad (3)\; (x_1 - x_2)^2$
\end{examplebox}

\textbf{解：}

由韦达定理：$x_1 + x_2 = 3$，$x_1 x_2 = 1$。

$(1)\; x_1^2 + x_2^2 = (x_1+x_2)^2 - 2x_1 x_2 = 9 - 2 = \boxed{7}$

$(2)\; \dfrac{1}{x_1} + \dfrac{1}{x_2} = \dfrac{x_1+x_2}{x_1 x_2} = \dfrac{3}{1} = \boxed{3}$

$(3)\; (x_1-x_2)^2 = (x_1+x_2)^2 - 4x_1 x_2 = 9 - 4 = \boxed{5}$

\begin{examplebox}[【例7】构造一元二次方程]
已知两个数的和为 $7$，积为 $12$，求这两个数。
\end{examplebox}

\textbf{解：}

设这两个数为 $x_1$、$x_2$，则 $x_1+x_2=7$，$x_1 x_2=12$。

由韦达定理的逆应用，$x_1$、$x_2$ 是方程 $x^2 - 7x + 12 = 0$ 的两个根。

分解因式：$(x-3)(x-4)=0$，所以 $x_1=3$，$x_2=4$。

这两个数为 $\boxed{3}$ 和 $\boxed{4}$。

\begin{examplebox}[【例8】含参数的方程讨论]
已知关于 $x$ 的方程 $x^2 - (2m+1)x + m^2 + m = 0$，证明：无论 $m$ 取何值，方程总有两个不相等的实数根。
\end{examplebox}

\textbf{证明：}

$\Delta = [-(2m+1)]^2 - 4(m^2+m)$

$= 4m^2 + 4m + 1 - 4m^2 - 4m$

$= 1 > 0$

因为 $\Delta = 1 > 0$ 恒成立，所以无论 $m$ 取何值，方程总有两个不相等的实数根。\qed

% ============================================================
% 第三部分：应用题
% ============================================================
\section*{\textcolor{apppurple}{三、应用题}}

\begin{appbox}[【应用1】增长率问题]
某工厂去年的产值为 $200$ 万元，计划在两年内使产值增加到 $288$ 万元。如果每年的增长率相同，求每年的增长率。
\end{appbox}

\textbf{解：}

设每年的增长率为 $x$（$x>0$）。

则今年产值为 $200(1+x)$ 万元，明年的产值为 $200(1+x)^2$ 万元。

由题意列方程：
\[
200(1+x)^2 = 288
\]
\[
(1+x)^2 = 1.44
\]
\[
1+x = \pm 1.2
\]
因为 $x>0$，取 $1+x=1.2$，得 $x=0.2=20\%$。

\textbf{答：}每年的增长率为 $\boxed{20\%}$。

\begin{appbox}[【应用2】面积问题]
用一根长为 $40$\,cm 的铁丝围成一个矩形框架，使矩形的面积为 $96$\,cm$^2$。求矩形的长和宽。
\end{appbox}

\textbf{解：}

设矩形的一边长为 $x$\,cm，则另一边长为 $\dfrac{40-2x}{2} = 20-x$（cm）。

由题意列方程：
\[
x(20-x) = 96
\]
\[
20x - x^2 = 96
\]
\[
x^2 - 20x + 96 = 0
\]
\[
(x-8)(x-12) = 0
\]
所以 $x_1 = 8$，$x_2 = 12$。

当 $x=8$ 时，另一边为 $12$；当 $x=12$ 时，另一边为 $8$。

\textbf{答：}矩形的长为 $\boxed{12}$\,cm，宽为 $\boxed{8}$\,cm。

\begin{appbox}[【应用3】利润问题]
某商店以每件 $30$ 元的价格购进一批商品。经市场调查发现，当售价为每件 $40$ 元时，每月可售出 $200$ 件；售价每增加 $1$ 元，月销量就减少 $10$ 件。若商店希望每月获得利润 $2160$ 元，那么每件商品的售价应定为多少元？
\end{appbox}

\textbf{解：}

设每件商品的售价比 $40$ 元增加 $x$ 元，则售价为 $(40+x)$ 元，每件利润为 $(40+x-30)=(10+x)$ 元，月销量为 $(200-10x)$ 件。

由题意列方程：
\[
(10+x)(200-10x) = 2160
\]
\[
2000 + 200x - 100x - 10x^2 = 2160
\]
\[
-10x^2 + 100x + 2000 = 2160
\]
\[
10x^2 - 100x + 160 = 0
\]
\[
x^2 - 10x + 16 = 0
\]
\[
(x-2)(x-8) = 0
\]
所以 $x_1=2$，$x_2=8$。

当 $x=2$ 时，售价为 $42$ 元；当 $x=8$ 时，售价为 $48$ 元。

\textbf{答：}每件商品的售价应定为 $\boxed{42}$ 元或 $\boxed{48}$ 元。

\begin{appbox}[【应用4】综合应用（面积问题变式）]
一块长方形草地的长比宽多 $5$\,m，要在草地中间开辟一条宽为 $1$\,m 的十字形道路（两条道路互相垂直），剩余草地的面积为 $150$\,m$^2$。求草地的长和宽。
\end{appbox}

\textbf{解：}

设草地的宽为 $x$\,m，则长为 $(x+5)$\,m。

十字形道路的面积为 $1\cdot x + 1\cdot(x+5) - 1\cdot1 = 2x + 4$（交叉部分减去重复计算的 $1$\,m$^2$）。

剩余草地面积 $= x(x+5) - (2x+4) = x^2 + 5x - 2x - 4 = x^2 + 3x - 4$。

由题意：
\[
x^2 + 3x - 4 = 150
\]
\[
x^2 + 3x - 154 = 0
\]
\[
(x-11)(x+14) = 0
\]
所以 $x_1 = 11$，$x_2 = -14$（不合题意，舍去）。

宽为 $11$\,m，长为 $11+5 = 16$\,m。

验证：草地面积 $11\times16=176$，道路面积 $2\times11+4=26$，剩余 $176-26=150$。\checkmark

\textbf{答：}草地的长为 $\boxed{16}$\,m，宽为 $\boxed{11}$\,m。

% ============================================================
% 第四部分：课堂练习 + 参考答案
% ============================================================
\section*{\textcolor{tipgreen}{四、课堂练习}}

\begin{enumerate}[leftmargin=*, label=\textbf{\arabic*.}]
  \item 用适当的方法解下列方程：
  \begin{enumerate}[label=(\arabic*)]
    \item $x^2 - 9 = 0$
    \item $x^2 + 6x + 5 = 0$
    \item $2x^2 - 3x - 2 = 0$
    \item $x^2 - 4x + 1 = 0$（用配方法）
  \end{enumerate}

  \item 判断下列方程根的情况（不解方程）：
  \begin{enumerate}[label=(\arabic*)]
    \item $x^2 + 2x + 5 = 0$
    \item $4x^2 - 12x + 9 = 0$
  \end{enumerate}

  \item 已知关于 $x$ 的方程 $x^2 + 2x + m - 1 = 0$ 有两个实数根，求 $m$ 的取值范围。

  \item 已知方程 $2x^2-5x+1=0$ 的两个根为 $x_1$、$x_2$，不解方程，求 $x_1^2+x_2^2$ 和 $\dfrac{1}{x_1}+\dfrac{1}{x_2}$ 的值。

  \item 某商品原价 $100$ 元，经过两次降价后售价为 $81$ 元。如果每次降价的百分率相同，求每次降价的百分率。
\end{enumerate}

\vspace{6pt}
\begin{knowledgebox}[参考答案]
\textbf{1.}\quad
(1)\; $x^2=9$，$x=\pm 3$，即 $x_1=3,\; x_2=-3$。

(2)\; $(x+1)(x+5)=0$，$x_1=-1,\; x_2=-5$。

(3)\; $(2x+1)(x-2)=0$，$x_1=-\dfrac{1}{2},\; x_2=2$。

(4)\; 配方：$x^2-4x+4=3$，$(x-2)^2=3$，$x=2\pm\sqrt{3}$。
即 $x_1=2+\sqrt{3},\; x_2=2-\sqrt{3}$。

\medskip
\textbf{2.}\quad
(1)\; $\Delta = 4-20 = -16 < 0$，方程没有实数根。

(2)\; $\Delta = 144-144 = 0$，方程有两个相等的实数根。

\medskip
\textbf{3.}\quad $\Delta = 4-4(m-1) = 8-4m \geqslant 0$，解得 $m \leqslant 2$。

\medskip
\textbf{4.}\quad 由韦达定理 $x_1+x_2=\dfrac{5}{2}$，$x_1 x_2=\dfrac{1}{2}$。

$x_1^2+x_2^2 = \left(\dfrac{5}{2}\right)^2 - 2\cdot\dfrac{1}{2} = \dfrac{25}{4}-1 = \dfrac{21}{4}$。

$\dfrac{1}{x_1}+\dfrac{1}{x_2} = \dfrac{x_1+x_2}{x_1 x_2} = \dfrac{5/2}{1/2} = 5$。

\medskip
\textbf{5.}\quad 设每次降价百分率为 $x$，则 $100(1-x)^2=81$，$(1-x)^2=0.81$，$1-x=0.9$（取正值），$x=0.1=10\%$。每次降价的百分率为 $10\%$。
\end{knowledgebox}

\vfill
\begin{center}
\small\textcolor{gray}{本讲义仅供教学参考，例题可根据学生水平调整难度。}
\end{center}

\end{document}
